\cotefile{11}
\scanpage{16}
\mspage{-2.29-}
% archive sheet number: 15

\uline{Proof}: Immediate, since $\overline{\lambda_1\unsure{\circ|\cdot}\lambda_2}=\bar{\lambda}_1+\bar{\lambda}_2$.

\uline{Corollary} \uline{2.16}: Let a reductive $G$ act on $X$, proper over $k$. Let $L\in\operatorname{Pic}^{G}(X)$ be ample, and let $x\in X_k$. Then $\{\delta\mid \unsure{\nu|v}^{L}(x;\delta)\leq 0\}$ is semi-convex, and $\{\delta\mid \nu^{L}(x,\delta)<0\}$ is convex.

\uline{Proof}: Immediate.

It follows that from this Corollary and from Theorem 2.1, that if \deleted{\unsure{$($}}$x\in X_k$ is not \uline{semi-stable}, then we can associate to $x$ the convex set $C$ of $\delta\in\Delta(G)$ for which $\nu^{L}(x,\delta)<0$. Roughly speaking, we can say that a parabolic subgroup $P$ is ``responsible'' for this breakdown of semi-stability if $C\cap\Delta_P(G)\neq\unsure{\phi|\emptyset}$. Since $C$ is convex, not too many parabolic subgroups $P$ are responsible; according to a conjecture of Tits, there is \inserted{\handwritten{even}} a natural way to find one $P$ which is \uline{most} responsible, i.e. $P(\delta)$ for $\delta=$ the ``center'' of $C$ (cf. Preface). Another consequence of Proposition 2.14 is:

\uline{Proposition} \uline{2.17}: Let a reductive $G$ act via $\sigma$ on $X$, proper over $k$. Let $L\in\operatorname{Pic}^{G}(x)$ be ample. Then there is a constant $K$ such that
\[
 \lvert\nu^{L}(x,\delta)\rvert\leq K
\]
for all $x\in X_k$ and $\delta\in\Delta(G)$. Moreover, for every subset $\handwritten{S\deleted{\unsure{\in}}\subset X}_k$, there is a $\delta_{\unsure{0|o}}\in\Delta(G)$ and $x_0\in S$ such that
\[
 \nu^{L}(x_0,\delta_0)\leq\nu^{L}(x,\delta)^{\unsure{\circ|.}}
\]
for every other $\delta\in\Delta(G)$, $x\in S$.
\scanpage{17}
% archive sheet number: 16

\inserted{provenant de} l'isom canonique entre $X\times Y$ et $Y\times X$. Soit alors
\[
 \varphi^{*}=s\varphi^{\circ}\colon V^{\circ}\longrightarrow V^{A}
\]
on constate alors que le composé
\[
 \psi\varphi^{*}=\psi(s\varphi^{\circ})=(\psi s)\varphi^{\circ}
 \colon V^{\circ}\longrightarrow k\text{-}\mathrm{mod}\,\mathrm{gr}
\]
% slanted
\inserted{\uline{\unsure{s'identifie|à vérifier}}} \deleted{est} \deleted{\unsure{canon|comme}} \uline{égal} à $\omega^{*}A^{\circ}\colon V^{\circ}\longrightarrow k\text{-}\mathrm{mod}\,\mathrm{gr}$.

\deleted{Ainsi, $V^{A}$ apparaît comme une espèce d'enveloppe}

\uline{Cas où $V$ a un objet final $e$}, \uline{et où} \uline{de plus} $A(e)=k$ \inserted{[Sinon, on remplacera $k$ par $A(e)$!]} Alors les $A(X)$ sont naturellement jaugés; si c'est des Alg. de Poincaré, il suffit donc de connaitre \deleted{\unsure{leurs|ces}} Alg. jaugées, $+$ les $f^{*}$ ($f\colon X\mathbin{-}Y$ dans $V$) pour \unsure{connaitre|connaître} $A$.

\[
\begin{aligned}
 \text{Soit}\qquad &u\in\operatorname{Hom}_{V^{A}}(X,Y)=A(X\times Y)\\
 &su\in\operatorname{Hom}_{V^{A}}(Y,X)=A(Y\times X)\\
 \text{donc}\qquad &\tilde u\colon A(X)\to\deleted{\unsure{X}}A(Y)
 \qquad \widetilde{su}\colon A(Y)\to A(X).
\end{aligned}
\]

Montrons que \inserted{si $u$ central dans $A(X\times Y)$, alors} $\tilde u$ et $\widetilde{su}$ sont transposés l'une de l'autre, pour les formes $k_{A(X)}$ et $k_{A(Y)}$ sur $A(X)$, $A(Y)$, i.e. on a pour $x\in A(X),\,y\in A(Y),\ u\in A(X\times Y)$
\[
\begin{aligned}
 k_{\deleted{\unsure{X}}Y}(y\deleted{\unsure{f}}u_{*}(x))&=k_X(\deleted{\unsure{f}}u^{*}(y)x)\qquad \text{i.e.}\\
 g_{*}(yu_{*}(x))&=f_{*}(u^{*}(y)x).
\end{aligned}
\]
(\uuline{Cas $u$ central} ?!) En effet

\diagram{\begin{tikzcd}[ampersand replacement=\&,column sep=large,row sep=small]
 \unsure{x} \& \unsure{y}\\
 A(\unsure{X|K})\arrow[r,no head,shift left,"u_*"]\arrow[d,"k_X"'] \& A(Y)\arrow[l,shift left,"u^*"']\arrow[d,"k_Y"]\\
 k\arrow[r,equal]\& k
\end{tikzcd}}

\diagram{\begin{tikzcd}[ampersand replacement=\&,column sep=small,row sep=small]
 \&X\times Y\arrow[dl,no head,"p_1"']\arrow[dr,no head,"p_2"]\&\\
 X\arrow[dr,no head,"f"']\&\&Y\arrow[dl,no head,"g"]\\
 \& e \&
\end{tikzcd}}

\[
\begin{aligned}
 g_*(yu_*(x))&=g_{\deleted{\unsure{*}}*}(\deleted{\unsure{u}}y p_{2*}(u p_1^*(x))=g_*p_{2*}(p_2^*(y)u p_1^*(x))=k(p_2^*(y)u p_1^{\unsure{*}}(\unsure{x}))\\
 f_*(u^*(y)x)&=f_*(p_{1*}(u p_2^*(y))x)=f_*p_{1*}(u p_2^*(y)p_1^*(x))=k(\deleted{\unsure{u}}u p_2^*(y)p_1^*(x))\\
 &=k(p_2^*(y)u p_1^*(x))\\[-2pt]
 &\unsure{>}\ \boxed{\text{mod. signe!}}
\end{aligned}
\]
\deleted{\illegible{1}}
\scanpage{18}
\mspage{-2.28-}
% archive sheet number: 17

\[
 L^{N}=\phi\unsure{\ast|{}^{*}}(\underline{\unsure{o|O}}(1))
\]
\[
 \text{Therefore:}\qquad \mu^{L}(x,\lambda)=\frac{1}{N}\mu^{o(1)}(\phi x,\lambda).
\]

As in \S1, the action of $G$ on $\handwritten{\mathbb{P}}_{n-1}$ is induced by a linear representation of $G$ in the affine cone $\handwritten{\mathbb{A}}^{n}$ over $\handwritten{\mathbb{P}}_{n-1}$. Suppose we choose coordinates in $\handwritten{\mathbb{A}}^{n}$ so that the representation of $T$ is diagonalized: i.e. $\alpha\in T_k$ acts via the matrix $\{\handwritten{\chi}_i(\alpha)\cdot\delta_{ij}\}$ for suitable characters $\handwritten{\chi}_1,\ldots,\handwritten{\chi}_n$ of $T$. Let $\handwritten{\chi}_i$ induce the linear functional $\bar{\handwritten{\chi}}_i$ on $\unsure{\Gamma|\Pi}^{\handwritten{\mathbb{R}}}(T)$, and as usual let a 1-PS $\lambda$ of $T$ induce the point $\bar{\lambda}\in\Gamma^{\handwritten{\mathbb{R}}}(T)$. Then if $(\phi x_1^*,\ldots,\phi x_n^*)$ are coordinates of a homogeneous point lying over $\phi x$, Proposition 2.3 asserts:
\[
 \mu^{o(1)}(\phi x,\lambda)=\max\bigl\{-\bar{\handwritten{\chi}}_i(\bar{\lambda})\bigm|i\text{ such that }\phi x_i^*\neq 0\bigr\}.
\]
We conclude:

\uline{Proposition} \uline{2.14}: Let a reductive $G$ act on $X$, proper over $k$. Let $T\subset G$ be a maximal torus, and let $L\in\operatorname{Pic}^{G}(X)$ be ample. Then there is a finite set of linear functionals $\ell_1,\ldots,\ell_n$ on $\Gamma^{\handwritten{\mathbb{R}}}(T)$ which are rational on $\Gamma^{\handwritten{\mathbb{Q}}}(T)$ with the following property:

\[
\begin{aligned}
&\text{for every }x\in X_k,\text{ there is a subset }I\subset\{1,\ldots,n\}\\
&\text{such that }\mu^{L}(x,\lambda)=\max\{\ell_i(\bar{\lambda})\mid i\in I\}\\
&\text{for all 1-PS's }\lambda\text{ of }T.
\end{aligned}
\]

\uline{Corollary} \uline{2.15}: If $\lambda_1,\lambda_2$ are 1-PS's of $T$ and $x\in X_k$, then
\[
 \mu^{L}(x,\lambda_{\unsure{1|J|j}}\unsure{\cdot|\circ}\lambda_2)\leq\mu^{L}(x,\lambda_1)+\mu^{L}(x,\lambda_2)
\]
\scanpage{19}
% archive sheet number: 18

\uline{Remarque} On n'a jamais, dans le calcul précédent, utilisé la structure additive des $A(X)$, seulement leur structure de monoïdes. Donc il y a lieu de formuler la théorie des correspondances en prenant \deleted{des} $X\mapsto A(X)$ \deleted{qui sont des} \inserted{foncteur dans les} ``monoïdes bivariants'' \ldots\,. \deleted{Il faudrait être \deleted{\unsure{précis|prudent}} \inserted{soigneux} pour les signes, dans le cas de monoïdes non commutatifs, en prenant des monoïdes gradués [les degrés en \unsure{seront additifs|seront astreints|anticommutatifs}]} Ce point de vue cependant ne semble pas très commode dans le cas où les 2 formules de projections \inserted{(gauche et droite)} ne sont pas valables l'une et l'autre, et qu'il faut introduire des signes. Il faudrait alors remplacer les monoïdes par des ``monoïdes jaugés'' $(A^i)_{i\in\unsure{\mathbb Z|Z|\mathbf{Z}}}$, loi de composition $A^i\times A^{\unsure{j|i}}\mathbin{\unsure{-|\to}} A^{i+j}$ (NB : L'ensemble somme des $A^i$ \deleted{\unsure{sera|est}} muni d'une structure de monoïde) et \deleted{\unsure{autom.} involutif d'ordre} \inserted{application bijective} $\sigma^i\colon A^i\mathbin{-} A^i$ satisfaisant $ {\sigma^i}^{2}=\mathrm{id}$, $\sigma^i=\mathrm{id}$ si $i$ pair, \deleted{pour} enfin chaque $A^*$ muni d'un \uline{degré} \inserted{$d(A^*)$} qui permet d'imposer le degré des $f_*$ des $(f_*,f^*)\colon A^{\unsure{*|\bullet}}\mathbin{-}B^*$ comme \deleted{\unsure{f}} $d(B^*)-d(A^{\unsure{*|\circ}})$. On peut alors \unsure{exprimer|imposer|imaginer} comment les signes doivent être calculés dans la formule de projection gauche (si on veut que la droite soit valable).
\scanpage{20}
\mspage{-2.26-}
% archive sheet number: 19

\uline{Proof}: $T_1$ and $T_2$ are both maximal tori in $P(\delta_1)\cap P(\delta_2)$ by lemma 2.9. Therefore they are conjugate in $P(\delta_1)\cap P(\delta_2)$. \hfill\uline{QED}

Now suppose $\delta_1,\delta_2$ are, in fact, in $\Delta(T_1)\cap\Delta(T_2)$: we must show that the lines joining them in $\Delta(T_1)$ and in $\Delta(T_2)$ are the same. Let $T_1=\gamma^{-1}\unsure{\cdot|\circ} T_2\cdot\gamma$ and suppose $\lambda_1,\lambda_2$ are 1-PS's of $T_2$ representing $\delta_1,\delta_2$, respectively. Then $\gamma^{-1}\cdot\lambda_1\cdot\gamma$, $\gamma^{-1}\cdot\lambda_2\cdot\gamma$ are 1-PS's of $T_1$ representing $\delta_1,\delta_2$ respectively. Therefore, we must show that the sets of points
\[
 \{\Delta(\lambda_1^{\handwritten{n_1}}\cdot\lambda_2^{\handwritten{n_2}})\}
 \qquad\text{and}\qquad
 \{\Delta\bigl((\gamma^{-1}\lambda_1\gamma)^{\handwritten{n_1}}\cdot(\gamma^{-1}\cdot\lambda_2\cdot\gamma)^{\handwritten{n_2}}\bigr)\}
\]
are the same. But this is clear provided that $\Delta(\lambda_1^{\handwritten{n_1}}\cdot\lambda_2^{\handwritten{n_2}})=\Delta(\gamma^{-1}\lambda_1^{\handwritten{n_1}}\cdot\lambda_2^{\handwritten{n_2}}\gamma)$; hence it follows it $P(\lambda_1)\cap P(\lambda_2)\subset P(\lambda_1^{\handwritten{n_1}}\cdot\lambda_2^{\handwritten{n_2}})$ - so that $\gamma\in P(\lambda_1^{\handwritten{n_1}}\cdot\lambda_2^{\handwritten{n_2}})$.

Therefore, we are reduced to proving that if $\lambda_1$, $\lambda_2$ are 1-PS's of $G$ which centralize each other, then
\begingroup\makeatletter\tagsleft@true\makeatother
\[
 P(\lambda_1\cdot\lambda_2)\supset P(\lambda_1)\cap P(\lambda_2) \tag*{(**)}
\]
\endgroup

To prove this, it is convenient to represent $G$ in $\mathrm{GL}(n)$ by some closed immersion $\phi\colon G\handwritten{\longrightarrow}\mathrm{GL}(n)$. Then, exactly as in the proof of Proposition 2.6, we have reduced the proof of (**) to the case $G=\mathrm{GL}(n)$. In this case, we can simultaneously diagonalize two commuting 1-PS's: say $\lambda_1(\alpha)$ is the matrix $\{\alpha^{\handwritten{r_i}}\cdot\delta_{ij}\}$, and $\lambda_2(\alpha)$ is the matrix $\{\alpha^{\handwritten{s_i}}\cdot\delta_{ij}\}$, for $\alpha\in(\handwritten{\mathbb{G}}_m)_k$. Then, as in Proposition 2.6,
\scanpage{21}
% archive sheet number: 20

Pour tout $X\in V$, soit $A'(X)\subset A(X)$, \inserted{avec $1_{A(X)}\in A'(X)$, stable par \unsure{isom|somme}.} Conditions équiv.

$1^\circ)$ Les $A'(X)$ stables par images inverses, images directes \deleted{par} tout morphisme de $V$, et par produits [ou encore, par produits $\boxtimes$].
% slanted
\margin{il faudrait donner un énoncé d'équivalence plus général, où on ne suppose \uline{pas} que \uline{\unsure{tt|H}} morphisme de $V$ \unsure{convient|une flèche}}

$2^\circ)$ \deleted{Soit} a) $u\in A'(X\times Y)$, $x\in A'(X)$, alors $u_*(x)\in A'(Y)$.

b) Si $x\in A'(X)$, $y\in A'(Y)$, alors $x\boxtimes y\in A'(X\times Y)$

c) Pour tout morphisme $X\mathbin{-}Y$ [il suffit de prendre les \unsure{.}] le graphe $\gamma_f=\Gamma_{f*}(1)\in A(X\times Y)$ est dans $A'(X\times Y)$.

\deleted{\uline{NB} Ces $A'$ correspondent aux ss-catégories (pas \unsure{forcément} pleines) de $V^A$, ayant mêmes objets que $V^A$, \unsure{et stables par}}

\uline{NB} Si $V$ a un objet final, les $A'$ correspondent exactement aux ss-catégories \deleted{de} $V_{A}$ qui ont \unsure{\={m} ens|tous ses|tjs ens|mêmes} d'objets que $A$, stables par $\boxtimes$ de flèches, et contenant les \deleted{\unsure{graphismes|graphes}} \deleted{\unsure{et leurs} \unsure{transposées}, les} morphismes de $V$ et leurs $s\Gamma_f$ (il suffit $s\delta_X$).

\uline{Dém} \quad $1^\circ)\Longrightarrow 2^\circ)$

$1^\circ$ forme \unsure{imprée|imprimée|primée|renforcée|simple} [avec $\boxtimes$] implique forme non imprée, car
% slanted
\margin{image \unsure{inverse|réciproque} par $\delta$}
\[
 x\otimes y=\delta^*(x\boxtimes y)\qquad[x,y\in A(X),\quad\delta=\mathrm{diag}_X\colon X\mathbin{-}X\times X]
\]

$1^\circ$ forme non imprée implique forme imprée, car
% slanted
\margin{images inverses par $p_i$, \deleted{\unsure{$p_2$}}}
\[
 x\boxtimes y=p_1^*(x)p_2^*(y).
\]

a) car $u_*(x)=p_{2*}(u p_1^*(x))$. [De $=$ pour $u^*(y)\in A'(X)$ si $y\in A'(Y)$]
% slanted
\margin{images directes, inverses par $\unsure{p_i|p_j|\mathrm{pr}\ \text{de }2}$}

b) clair

c) clair, car $1_X\in A'(X)$ par hyp.
% slanted
\margin{image directe \inserted{de $1$} par un graphe}
% slanted
\margin{$+$ produits}

$2^\circ)\Longrightarrow 1^\circ)$. Stabilité par images \deleted{directes} directes : car $f_*(x)=\deleted{\unsure{\Gamma_f|\gamma_f}}\gamma_{f*}(x)$, on applique a) et c).

Stabilité par images \deleted{directes} \inserted{inverses} : Car $f^*(y)=({}^{t}\gamma_f)_*(y)$, et comme $\gamma_f\in A'(X\times Y)$, on a $ {}^{t}\gamma_f\in A'(Y\times X)$ (stabilité par isom),

Stabilité par $\boxtimes$ claire.
\scanpage{22}
\mspage{-2.24-}
% archive sheet number: 21

of the terminology ``flag complex''.

A second direction in which we can make explicit this structure on $\Delta(G)$ is to show that the skeletons are pasted together in a sufficiently nice way so that various constructions which can be carried out in each skeleton, can also be carried out in $\Delta(G)$\unsure{.|:|;} The principal one we have in mind is drawing a great circle through two points. One approach to this is to complete the complex $\Delta(G)$ with respect to the metric $\handwritten{\rho}$ mentioned above: then Tits has shown that the geodesic segments on this completion are precisely the segments of great circles of length $\leq\pi$, on the true spheres obtained by completing the skeletons $\Delta(T)$. Since we are avoiding this metric, we define these lines, and other auxiliary concepts, directly:

\uline{Definition} \uline{2.8}: A pair of points $\delta_1,\delta_2\in\Delta(T)$ is \uline{antipodal} if there is a 1-PS $\lambda$ of $G$ such that $\delta_1=\handwritten{\Delta}(\handwritten{\lambda})$ and $\delta_2=\Delta(\lambda^{-1})$.

If $G$ is a torus, then every point of $\Delta(G)$ has a unique antipodal point, i.e. its antipodal on the sphere is the usual sense. More general, if $\delta_1$ is in the skeleton $\Delta(T)$ of $\Delta(G)$, for some maximal torus $T$, then the antipodal $\delta_2$ of $\delta_1$ as a point of $\Delta(T)$, is one of its antipodal points in $\Delta(G)$. In fact, these are the only points antipodal to $\delta_1$:

\uline{lemma} \uline{2.12}: If $\delta_1,\delta_2$ are two points in the skeleton $\Delta(T)$ of $\Delta(G)$, for some maximal torus $T$, then $\delta_1$ and $\delta_2$ are antipodal as points of $\Delta(G)$ if and only if they are antipodal as points of the sphere\unsure{'} $\Delta(T)$.
\scanpage{23}
% archive sheet number: 22

\deleted{\unsure{Intersection} \unsure{Sous-an|Sommes}}

Sous-Anneau bivariant sur $V$. [Satisfait à son compte les form. de proj. et d'échange] Intersection de tels est idem; d'où notion de foncteur ``stable'' \uline{engendré}. Soit $A_{\mathrm{sp}}\deleted{\unsure{\in}}\subset A$ le plus petit ss-foncteur stable engendré.

\uline{NB} $A_{\mathrm{sp}}(X)$ contient les $f_*(1_Y)$ pour $f\colon Y\mathbin{\unsure{-|\to}}X$ dans $V$. Dans beaucoup de cas importants, il est $=$ \unsure{identiques|identique} à l'ens des $f_*(1_Y)$ pour tous les $f\colon Y\mathbin{-}X$ dans $\mathrm{Ob}(V_{/X})$. Cela \unsure{marche|vaudra|voudra} p. ex. si $V$ la catégorie des schémas projectifs lisses \unsure{équidim} sur $k$ alg. clos, dans les cas envisagés plus haut, - du moins si on admet la résolution des \uline{singularités}.

% slanted
\margin{$A'$ \inserted{$\subset A$} donné, \uline{stable}}
Une application $A(X)\times\cdots\times A(X_n)\mathbin{-}A(X)$ est dite \uline{de type} $A'$ si elle est de la forme
\[
 x_1,\ldots,x_n\longmapsto u(x_1\boxtimes\cdots\boxtimes x_n)
\]
avec $u\in A'(X_1\times\cdots\times X_n\times X)$. En \deleted{\unsure{un cas}} particulier, ceci est défini si $n=1$. Composant une \inserted{telle} application $\varphi$ avec une application $\psi\colon A(X)\mathbin{-}A(X')$ qui est de type $A'$, on trouve une application qui est de type $A'$.

\deleted{\unsure{De} plus \unsure{génér.|général.} De \unsure{m.}}, \inserted{Plus gén.} si on compose des applications \deleted{Donnons-nous des \unsure{fonctions|applications} $(\unsure{\psi|X}_{\unsure{i|h}})_{\unsure{h|i}\in I}$}
\[
\begin{aligned}
 \varphi_{\unsure{i|1}}(x_1,\ldots,x_n)&\colon A(X_1)\times\cdots\times A(X_n)\mathbin{-}A(X)\\
 \psi(y_1,\ldots,y_m)&\colon A(Y_1)\mathbin{\unsure{-|\times}}\cdots\times A(Y_m)\mathbin{-}A(Y)\\
 \mathbin{\unsure{-|\to}}\ \text{avec}\qquad\chi&\colon A(X)\times A(Y)\longrightarrow A(Z)
\end{aligned}
\]
\scanpage{24}
\mspage{-2.23-}
% archive sheet number: 23

Therefore, in view of lemma 2.9, $\Delta(G)$ can be viewed as the result of pasting together spheres, one for each $T\subset G$, with sufficient identifications so that any two points are both on at least one sphere. The subsets $\Delta(T)$ will be called \uline{skeletons} of $\Delta(G)$. We can push this structure further in 2 ways: first the spheres $\Delta(T)$ can be broken up by considering the function $P(\delta)$:

\uline{Definition} \uline{2.7}: For all parabolic subgroups $P$, let $\Delta_{\handwritten{\unsure{P|p}}}(G)$ be the set of $\delta\in\Delta(G)$ such that $P(\delta)\supset P$.

To describe $\Delta_P(G)$ explicitly, let $T$ be a maximal torus of $P$. Then $\Delta_P(G)\subset\Delta(T)$ by lemma 2.9. Suppose $P$ is spanned by $T$ and by additive subgroups $\handwritten{\mathbb{G}}_a\subset P$ such that
\[
 \alpha^{-1}\cdot\gamma\cdot\alpha=\handwritten{\chi}(\alpha)\cdot\gamma
\]
for $\alpha\in T_k$, $\gamma\in(\handwritten{\mathbb{G}}_a)_k$, and a suitable character $\handwritten{\chi}$ of $T$. Then one checks that if $\delta\in\Delta(T)$ is induced by $\bar{\lambda}\in\unsure{\Gamma|\Pi}^{\handwritten{\mathbb{R}}}(T)$, and if $\handwritten{\chi}$ induces the linear functional $\bar{\handwritten{\chi}}$ on $\Gamma^{\handwritten{\mathbb{R}}}(T)$, then
\[
 P(\delta)\supset P\quad\handwritten{\Longleftrightarrow}\quad\handwritten{\forall}\text{ such }\handwritten{\chi},\quad\bar{\handwritten{\chi}}(\bar{\lambda})\leq 0.
\]
Therefore, $\Delta_P(G)$ is simple the set of rational points on an intersection of hemi-spheres, (each determined by a rational linear functional, in fact). If $G$ is actually semi-simple, it can be shown that $\Delta_{\handwritten{P}}(G)$ is even the set of rational points on a spherical simplex, and hence the collection of $\Delta_{\handwritten{P}}(G)$ constitutes a ``triangulation'', and $\Delta(G)$ is a complex formed from the simplices $\Delta_{\handwritten{P}}(G)$. This is the motivation
\scanpage{25}
% archive sheet number: 24

avec $\varphi,\psi,\chi$ de type $A'$, on trouve une application \deleted{$A(X)$} $\chi(\varphi,\psi)\colon A(X_1)\times\cdots\times A(X_n)\times A(Y_1)\times\cdots\times A(Y_m)\longrightarrow A(Z)$ qui est de type $A'$.

\margin{Donc le produit de deux applications \deleted{$A'$-\unsure{admissibles}} \inserted{de type $A'$} [de $=$ orig. et but] est \deleted{$A'$-\unsure{adm}} de type $A'$; de $=$ le produit par un élément de $A'(X)$.}

\uline{Application} \deleted{Soit $E\subset V$, et pour} Pour tt $X\in V$, supposons donné $E_X\in A(X)$ (qui peut être vide). Alors le plus petit syst. stable $B\subset A$ contenant $A'$ et les $E_X$ est défini ainsi :
\[
\begin{aligned}
B(X)={}&\deleted{\text{\unsure{ens}}}\inserted{\deleted{\text{\illegible{1}}} A'(X)\cup}\text{ \unsure{des|ens} éléments \inserted{de \deleted{$A(X)$ \unsure{de|dans} $A'(X)$.}} de la forme }\inserted{1\text{ ou }}\varphi(x_1,\ldots,\unsure{x_n|a_n})\text{ où}\\
&\varphi\text{ est une application } A'\text{-\unsure{admissibles|admissible} de } A(X_1)\times\cdots\times A(X_n)\\
&\text{dans }A(X),\quad x_i\in E(X_i)\text{ pour }i=1,-n,\text{ les }X_i\in\mathrm{Ob}\,V.
\end{aligned}
\]
En effet, les $B(X)$ \inserted{ainsi définis} sont stables par images directes et inverses, \deleted{il \unsure{reste}} par morphismes de $V$ \inserted{et \unsure{contient|contiennent} les $1$, et les $A'$;} reste à prouver qu'ils sont stables par multiplication. OK.
\scanpage{26}
\mspage{-2.22-}
% archive sheet number: 25

of some torus $T\subset H$. If $\epsilon_1\neq\epsilon_2$, then there do not exist positive integers $n$ and $m$ such that $\lambda_1(\handwritten{\alpha^n})=\lambda_2(\handwritten{\alpha^m})$, all $\alpha\in\handwritten{\mathbb{G}}_m$. But then there is a character $\handwritten{\chi}$ of $T$ such that $\handwritten{\chi}\circ\lambda_1=0$, (the trivial homomorphism from $\handwritten{\mathbb{G}}_m$ to $\handwritten{\mathbb{G}}_m$), and $\handwritten{\chi}\circ\lambda_2\neq0$. Since the kernel of $\phi$ is finite, there is an integer $N$ such that $\handwritten{\chi}^N$ is induced by a character $\handwritten{\chi'}$ of $\phi(T)$. Now suppose that $\phi_*\epsilon_1=\phi_*\epsilon_2$. Then $P(\phi\circ\lambda_1)=P(\phi_*\epsilon_1)=P(\phi_*\epsilon_2)=P(\phi\circ\lambda_2)$. Let $T^{\handwritten{\prime\prime}}$ be the intersection of $\phi(T)$, with the radical of $P(\phi\circ\lambda_1)$, and let $\handwritten{\chi''}$ be the restriction of $\handwritten{\chi'}$ to $T$. According to the structure theory of parabolic subgroups (cf.\hspace{4em}), some power of $\handwritten{\chi''}$ extends to a character $\handwritten{\chi'''}$ of $P(\phi\circ\lambda_1)$. But since $\phi\circ\lambda_1$ and $\phi\circ\lambda_2$ are 1-parameter subgroups of $P(\phi\circ\lambda_1)$, hence of $T^{\handwritten{\prime\prime}}$, we have $\handwritten{\chi'''}\circ(\phi\circ\lambda_1)=0$, and $\handwritten{\chi'''}\circ(\phi\circ\lambda_2)\neq0$. But then, referring to the definition of $\Delta(G)$, one sees that this precludes the possibility $\Delta(\phi\circ\lambda_1)=\Delta(\phi\circ\lambda_2)$. Therefore $\phi_*\epsilon_1\neq\phi_*\epsilon_2$. \uline{QED}

\uline{Corollary} \uline{2.11}: If $\phi$ is an isogeny, $\phi_*$ is an isomorphism.

\uline{Proof}: It remains to show that $\phi_*$ is surjective. But if $\lambda$ is any 1-parameter subgroup of $G$, then for some positive integer $n$, and for some 1-parameter subgroup $\mu$ of $H$, $\phi\circ\mu(\alpha)=\lambda(\alpha^{\handwritten{\unsure{h|n}}})$, $\alpha\in\handwritten{\mathbb{G}}_m$. \uline{QED}

Suppose we apply Proposition 2.1\deleted{\unsure{0|1}}\handwritten{\unsure{1|0}} to the inclusion of a maximal torus $T$ in $G$. This induces an inclusion of $\Delta(T)$ in $\Delta(G)$. But $\Delta(T)$ is very simple - it is essentially the set of rational points on a sphere. In fact, it follows immediately from the definition that $\Delta(T)$ is the set of rays in $\unsure{\Gamma|\Pi}^{\handwritten{\mathbb{R}}}(T)$ which contain a point of $\Gamma^{\handwritten{\mathbb{Q}}}(T)$.
\scanpage{27}
% archive sheet number: 26

\deleted{
\diagram{\begin{tikzcd}[ampersand replacement=\&]
 A\times B\arrow[r,"s_{A,B}"]\arrow[d,"f\times g"] \& B\times A\arrow[d,"g\times f"]\\
 A'\times B'\arrow[r,"s_{A',B'}"] \& B'\times A'
\end{tikzcd}}
\[
\begin{aligned}
(g\times f)_*(s(x\otimes y))&=(g\times f)_*(y\otimes x)(-1)^{xy}
=\unsure{\{}g_*(y)\otimes f_*(\unsure{y|x})(-1)^{xy+\text{\textcircled{$fx$}}}\\
s(f\times g)_*(x\otimes y)&=s(f_*(x)\otimes g_*(y))(-1)^{\text{\textcircled{$gx$}}}
=g_*(y)\otimes f_*(\unsure{x})(-1)^{\text{\textcircled{$gx$}}+(x+f)(y+g)}
\end{aligned}
\]
\[
\begin{aligned}
&(-1)^{\deleted{xy}+fx+\deleted{gx}+\deleted{gx}+fy+\deleted{xy}+fg}
\qquad\deleted{=(-1)^{(f+x)(g+y)}}\\
&\unsure{\shortparallel}\qquad\\[-.4em]
&(-1)^{f(x+y+g)}\qquad(-1)^{\deleted{\unsure{xy+}}fy+gx+fg}
\end{aligned}
\]
}
\noindent\rule{.91\linewidth}{.4pt}

\uline{A}

1) Catégorie \inserted{$V^\circ$} des alg. de Poincaré $A$ sur $k$ : alg. gr. \unsure{comm.} proj. de t. f. à degrés $\geq0$, munie d'un hom. lin. $k_A\colon A\mathbin{-}k$, de degré $d^\circ(A)$ dépendant de $A$ tel que $(x,y)\longmapsto k_A(xy)$ soit une forme \inserted{strictement} non dégénérée. \deleted{\unsure{pour} \unsure{strict}}

Les hom. \inserted{de $V^\circ$} : ce sont les hom. d'algèbres gr. \unsure{degrés|unités}, sans plus (donc il n'est pas question dans la définition de $V$ hom des $k_A$)

\margin{$\unsure{Z|2}$}
\uline{NB} Un hom de cette catégorie n'est pas néc. compatible avec \inserted{les} $k_A,k_B$

2) Foncteur contravariant $\varphi\colon (V^\circ)^\circ=V\longrightarrow\deleted{\unsure{M}}\unsure{\mathrm{Mod\,gr.d}|\mathrm{Modgrad}}(k)$

a) \unsure{Sur|sur} les objets : $\varphi(A)=A$ où on oublie la multiplication

b) Sur les morphismes : soit $f^*\colon B\mathbin{-}A$, on pose $\varphi(f^*)=f_*\colon A\longrightarrow B$ défini par la condition que
\[
k_A(f^*(b)a)=k_B(bf_*(a))
\]
On vérifie $\varphi(g^*f^*)=\varphi(\unsure{f^*|f})\varphi(g^*)$, $\varphi(\mathrm{id})=\mathrm{id}$, et formule de proj.
\[
f_*(a)b=f_*(af^*b)
\]
enfin \deleted{$k_B(f_*(a))$} $k_B f_*=k_A$
\scanpage{28}
\mspage{-2.14-}
% archive sheet number: 27

is still closed. Therefore, by definition, $x$ is a point of $\handwritten{U}$.

Therefore, a quasi-projective geometric quotient $X/\handwritten{\mathbb{G}}_m$ exists. But this implies that the action of $\handwritten{\mathbb{G}}_m$ on $X$ is separated. Therefore, by lemma 0.5, the action of $\handwritten{\mathbb{G}}_m$ on $X$ is proper. \uline{QED}

\fbox{\S2.} In the last section, we have related the concept of stability with respect to actions of $G$ to the function $\mu^L(x;\lambda)$ involving a 1-parameter subgroup $\lambda$ of $G$. In this section,we shall consider the dependence of $\mu$ on $\lambda$. The analysis leads naturally to an ungainly but remarkable metric space investigated extensively by J. Tits, which we will call the \uline{flag} \uline{complex} of $G$. The first step is:

\uline{Definition} \uline{2.3}/\uline{Proposition} \uline{2.6}: Let $G$ be a reductive algebraic group and l\deleted{\unsure{e}}\handwritten{e}t $\lambda$ be a 1-PS of $G$. Then there is a unique algebraic subgroup $P(\lambda)\subset G$ such that:
\[
\begin{aligned}
\gamma\in P(\lambda)_k\quad\handwritten{\Longleftrightarrow}\quad\lambda(\alpha)\cdot\gamma\cdot\lambda(\alpha^{-1})&\text{ has a specialization in }G\\
&\text{ when }\alpha\in(\handwritten{\mathbb{G}}_m)_k\text{ specializes to }0.
\end{aligned}
\]
Moreover, $P(\lambda)$ is a parabolic subgroup of $G$, $\lambda$ is a 1-PS of the radical of $P(\lambda)$, and the specialization $\gamma'$ of $\lambda(\alpha)\cdot\gamma\cdot\lambda(\alpha^{-1})$, as $\alpha\handwritten{\longrightarrow}0$, centralizes $\lambda$.

\uline{Proof}: Most of these facts can be seen most easily by means of a faithful representation $\phi\colon G\handwritten{\longrightarrow}\mathrm{GL}(n)$. Recall that the image $\phi(G)$ is \uline{closed} in $\mathrm{GL}(n)$. Therefore, $P(\lambda)=\phi^{-1}(P(\phi\circ\lambda))$, at least as a set of closed points. Now by composing $\phi$ \deleted{\unsure{by}} \inserted{\handwritten{with}} a suitable inner automorphism of $\mathrm{GL}(n)$, we may assume that $\phi\circ\lambda$ is diagonalized, i.e.
